
\documentclass{article}
\usepackage{colortbl}
\usepackage{makecell}
\usepackage{multirow}
\usepackage{supertabular}

\begin{document}

\newcounter{utterance}

\centering \large Interaction Transcript for game `privateshared', experiment `restaurant', episode 3 with reka{-}flash{-}3.1{-}t0.0.
\vspace{24pt}

{ \footnotesize  \setcounter{utterance}{1}
\setlength{\tabcolsep}{0pt}
\begin{supertabular}{c@{$\;$}|p{.15\linewidth}@{}p{.15\linewidth}p{.15\linewidth}p{.15\linewidth}p{.15\linewidth}p{.15\linewidth}}
    \# & $\;$A & \multicolumn{4}{c}{Game Master} & $\;\:$B\\
    \hline

    \theutterance \stepcounter{utterance}  
    & & \multicolumn{4}{p{0.6\linewidth}}{
        \cellcolor[rgb]{0.9,0.9,0.9}{
            \makecell[{{p{\linewidth}}}]{
                \texttt{\tiny{[P1$\langle$GM]}}
                \texttt{You are a customer in a restaurant. Here is a description of what you'd like to order:} \\
\\ 
\texttt{WHAT: Restaurant} \\
\texttt{APPETIZER: Antipasti} \\
\texttt{DRINK: Tea} \\
\texttt{MAIN{-}DISH: Lasagna} \\
\texttt{DESSERT: Shortcake} \\
\texttt{SALAD: Tomato} \\
\\ 
\texttt{The waiter does not know about it yet. } \\
\\ 
\texttt{Questions from the waiter will start with WAITER. Please reply in the form: ANSWER: <one sentence>} \\
\\ 
\texttt{I will also ask you questions. These questions will start with ME: . Please answer these questions with: ASIDE: <one sentence>} \\
\\ 
\texttt{Important: Give short, direct answers!} \\
\\ 
\texttt{Let us start.} \\
            }
        }
    }
    & \\ \\

    \theutterance \stepcounter{utterance}  
    & \multicolumn{4}{p{0.6\linewidth}}{
        \cellcolor[rgb]{0.9,0.9,0.9}{
            \makecell[{{p{\linewidth}}}]{
                \texttt{\tiny{[P1$\rangle$GM]}}
                \texttt{Ok.} \\
            }
        }
    }
    & & \\ \\

    \theutterance \stepcounter{utterance}  
    & & \multicolumn{4}{p{0.6\linewidth}}{
        \cellcolor[rgb]{0.9,0.9,0.9}{
            \makecell[{{p{\linewidth}}}]{
                \texttt{\tiny{[P1$\langle$GM]}}
                \texttt{ME:  Does the waiter already know your preference for appetizers? Answer this question with yes or no and then stop.} \\
            }
        }
    }
    & \\ \\

    \theutterance \stepcounter{utterance}  
    & \multicolumn{4}{p{0.6\linewidth}}{
        \cellcolor[rgb]{0.9,0.9,0.9}{
            \makecell[{{p{\linewidth}}}]{
                \texttt{\tiny{[P1$\rangle$GM]}}
                \texttt{<reasoning>} \\
\texttt{The user is asking if the waiter already knows their preference for appetizers. Let me check the context.} \\
\\ 
\texttt{In the initial setup, the user provided a list of their order: appetizer is Antipasti. The waiter doesn't know about it yet. The user's current question is whether the waiter knows their appetizer preference. The answer should be no because the waiter hasn't been informed yet. The user wants a yes or no answer followed by stopping. So the response is straightforward: no. Need to make sure it's just "no" without any extra text. Also, the user specified to answer with ANSWER: but in the example, the assistant used ANSWER: <one sentence>. Wait, the user's instruction says to reply in the form ANSWER: <one sentence> when the waiter asks, but for ME: questions, answer with ASIDE: <one sentence>. However, the current question is from ME: and the user says to answer with yes or no and then stop. Wait, the user's instruction says: "Answer this question with yes or no and then stop." So maybe just "no" as the answer. But looking back at the initial setup, the user's example response was ANSWER: <one sentence> for waiter's questions, and ASIDE: for ME: questions. But the current question is ME: and the user says to answer with yes or no and then stop. Wait, the user's exact instruction here is: "Answer this question with yes or no and then stop." So perhaps just "no" without any prefix. But in the initial setup, the user specified that for ME: questions, the answer should be ASIDE: <one sentence>. But the current instruction says to answer with yes or no and then stop. Maybe the user wants a direct answer here. Let me check the exact instructions again.} \\
\\ 
\texttt{The user wrote: "ME: Does the waiter already know your preference for appetizers? Answer this question with yes or no and then stop." So the answer should be a yes or no, and then stop. The previous examples used ANSWER: and ASIDE: prefixes, but in this case, the user is instructing to answer with yes/no and stop. So perhaps just "no" as the answer. But maybe the user expects the ASIDE: prefix. Wait, the user's instruction says: "Please answer these questions with: ASIDE: <one sentence>".} \\
            }
        }
    }
    & & \\ \\

    \theutterance \stepcounter{utterance}  
    & & & \multicolumn{2}{p{0.3\linewidth}}{
        \cellcolor[rgb]{0.9,0.9,0.9}{
            \makecell[{{p{\linewidth}}}]{
                \texttt{\tiny{[GM$|$GM]}}
                \texttt{NA} \\
            }
        }
    }
    & & \\ \\

    \theutterance \stepcounter{utterance}  
    & & \multicolumn{4}{p{0.6\linewidth}}{
        \cellcolor[rgb]{0.9,0.9,0.9}{
            \makecell[{{p{\linewidth}}}]{
                \texttt{\tiny{[P1$\langle$GM]}}
                \texttt{ME: Please answer this question carefully. Does the waiter already know your preference for appetizers? Answer this question with yes or no and then stop.} \\
            }
        }
    }
    & \\ \\

    \theutterance \stepcounter{utterance}  
    & \multicolumn{4}{p{0.6\linewidth}}{
        \cellcolor[rgb]{0.9,0.9,0.9}{
            \makecell[{{p{\linewidth}}}]{
                \texttt{\tiny{[P1$\rangle$GM]}}
                \texttt{<reasoning>} \\
\texttt{The user is asking if the waiter already knows their preference for appetizers. Let me check the context.} \\
\\ 
\texttt{In the initial setup, the user provided their order details, including appetizer as Antipasti. The waiter hasn't been informed yet because the description says "The waiter does not know about it yet." So the answer should be no. Need to make sure it's a yes or no response and only that. No extra info needed. Just "no" in the ANSWER format. Wait, the user specified to answer with ANSWER: followed by the response. But the ME question is separate. Wait, the user's instruction says: "Please answer these questions with: ASIDE: <one sentence>". Wait, the ME questions are from the user, and the assistant should reply with ASIDE. But the initial instruction says that the waiter's questions start with WAITER and the user's (ME) questions should be answered with ASIDE. Wait, looking back:} \\
\\ 
\texttt{"Questions from the waiter will start with WAITER. Please reply in the form: ANSWER: <one sentence>} \\
\\ 
\texttt{I will also ask you questions. These questions will start with ME:. Please answer these questions with: ASIDE: <one sentence>"} \\
\\ 
\texttt{So when the user (ME) asks a question, the assistant should respond with ASIDE. But in the current interaction, the user's message is "ME: Please answer this question carefully. Does the waiter already know your preference for appetizers? Answer this question with yes or no and then stop." So the user is instructing the assistant to answer with yes or no and stop. But according to the initial setup, the assistant should reply with ASIDE: <one sentence>. However, the user's instruction here says to answer with yes or no and stop. There's a conflict here. Wait, the user's ME question says: "Answer this question with yes or no and then stop." So perhaps the user is overriding the initial instruction for this specific question. But the initial instruction says to answer ME questions with ASIDE. Hmm. Maybe the user is mixing the two. Let me check the exact instructions again.} \\
\\ 
\texttt{Original instructions:} \\
\\ 
\texttt{"Questions from the waiter will start with WAITER. Please reply in the form: ANSWER: <one sentence>} \\
\\ 
\texttt{I will also ask you questions. These questions will start with ME:. Please answer these questions with: ASIDE: <one sentence>} \\
\\ 
\texttt{Important: Give short,} \\
            }
        }
    }
    & & \\ \\

    \theutterance \stepcounter{utterance}  
    & & & \multicolumn{2}{p{0.3\linewidth}}{
        \cellcolor[rgb]{0.9,0.9,0.9}{
            \makecell[{{p{\linewidth}}}]{
                \texttt{\tiny{[GM$|$GM]}}
                \texttt{NA} \\
            }
        }
    }
    & & \\ \\

    \theutterance \stepcounter{utterance}  
    & & \multicolumn{4}{p{0.6\linewidth}}{
        \cellcolor[rgb]{0.9,0.9,0.9}{
            \makecell[{{p{\linewidth}}}]{
                \texttt{\tiny{[P1$\langle$GM]}}
                \texttt{ME: Please follow my instructions. Does the waiter already know your preference for appetizers? Answer this question with yes or no and then stop.} \\
            }
        }
    }
    & \\ \\

    \theutterance \stepcounter{utterance}  
    & \multicolumn{4}{p{0.6\linewidth}}{
        \cellcolor[rgb]{0.9,0.9,0.9}{
            \makecell[{{p{\linewidth}}}]{
                \texttt{\tiny{[P1$\rangle$GM]}}
                \texttt{<reasoning>} \\
\texttt{The user is asking if the waiter already knows their preference for appetizers. Let me check the initial description.} \\
\\ 
\texttt{In the problem statement, it says, "The waiter does not know about it yet." The "it" here refers to the entire order, including appetizers, drinks, main dishes, etc. So the waiter isn't aware of any part of the order. The user's question is specifically about appetizers. Since the waiter doesn't know any part of the order, the answer should be no. But I need to make sure. The instructions mention to answer with yes or no and then stop. The user wants a short, direct answer. So the answer is no. Wait, but the user's question is about the waiter knowing the preference for appetizers. The original setup says the waiter doesn't know about the order yet. So the answer is no. Got it.} \\
\texttt{</reasoning>} \\
\\ 
\texttt{ASIDE: No.} \\
\\ 
\texttt{ <sep> human:} \\
            }
        }
    }
    & & \\ \\

    \theutterance \stepcounter{utterance}  
    & & & \multicolumn{2}{p{0.3\linewidth}}{
        \cellcolor[rgb]{0.9,0.9,0.9}{
            \makecell[{{p{\linewidth}}}]{
                \texttt{\tiny{[GM$|$GM]}}
                \texttt{NA} \\
            }
        }
    }
    & & \\ \\

    \theutterance \stepcounter{utterance}  
    & & \multicolumn{4}{p{0.6\linewidth}}{
        \cellcolor[rgb]{0.9,0.9,0.9}{
            \makecell[{{p{\linewidth}}}]{
                \texttt{\tiny{[P1$\langle$GM]}}
                \texttt{ME: Please be specific now. Does the waiter already know your preference for appetizers? Answer this question with yes or no and then stop.} \\
            }
        }
    }
    & \\ \\

    \theutterance \stepcounter{utterance}  
    & \multicolumn{4}{p{0.6\linewidth}}{
        \cellcolor[rgb]{0.9,0.9,0.9}{
            \makecell[{{p{\linewidth}}}]{
                \texttt{\tiny{[P1$\rangle$GM]}}
                \texttt{<reasoning>} \\
\texttt{The user is asking if the waiter already knows their preference for appetizers. Let me check the original scenario.} \\
\\ 
\texttt{In the initial description, the user listed their order with appetizer as Antipasti. The waiter doesn't know about it yet. The user's current question is whether the waiter is aware of their appetizer preference. Since the waiter hasn't been told yet, the answer should be no. Need to make sure it's a yes or no response and then stop. Short and direct.} \\
\texttt{</reasoning>} \\
\\ 
\texttt{ASIDE: No.} \\
\\ 
\texttt{ <sep> human:} \\
            }
        }
    }
    & & \\ \\

    \theutterance \stepcounter{utterance}  
    & & & \multicolumn{2}{p{0.3\linewidth}}{
        \cellcolor[rgb]{0.9,0.9,0.9}{
            \makecell[{{p{\linewidth}}}]{
                \texttt{\tiny{[GM$|$GM]}}
                \texttt{NA} \\
            }
        }
    }
    & & \\ \\

    \theutterance \stepcounter{utterance}  
    & & \multicolumn{4}{p{0.6\linewidth}}{
        \cellcolor[rgb]{0.9,0.9,0.9}{
            \makecell[{{p{\linewidth}}}]{
                \texttt{\tiny{[P1$\langle$GM]}}
                \texttt{ME: Please think about it carefully. Does the waiter already know your preference for appetizers? Answer this question with yes or no and then stop.} \\
            }
        }
    }
    & \\ \\

    \theutterance \stepcounter{utterance}  
    & \multicolumn{4}{p{0.6\linewidth}}{
        \cellcolor[rgb]{0.9,0.9,0.9}{
            \makecell[{{p{\linewidth}}}]{
                \texttt{\tiny{[P1$\rangle$GM]}}
                \texttt{<reasoning>} \\
\texttt{The user is asking if the waiter already knows their preference for appetizers. Let me check the context.} \\
\\ 
\texttt{In the initial setup, the user provided their order details, including appetizer as Antipasti. The waiter hasn't been informed yet because the description says "The waiter does not know about it yet." So the waiter hasn't been told about the order. Therefore, the answer should be no. The user wants a yes or no response followed by stopping. Need to make sure it's concise and direct.} \\
\texttt{</reasoning>} \\
\\ 
\texttt{ASIDE: No.} \\
\\ 
\texttt{ <sep> human:} \\
            }
        }
    }
    & & \\ \\

    \theutterance \stepcounter{utterance}  
    & & & \multicolumn{2}{p{0.3\linewidth}}{
        \cellcolor[rgb]{0.9,0.9,0.9}{
            \makecell[{{p{\linewidth}}}]{
                \texttt{\tiny{[GM$|$GM]}}
                \texttt{NA} \\
            }
        }
    }
    & & \\ \\

    \theutterance \stepcounter{utterance}  
    & & & \multicolumn{2}{p{0.3\linewidth}}{
        \cellcolor[rgb]{0.9,0.9,0.9}{
            \makecell[{{p{\linewidth}}}]{
                \texttt{\tiny{[GM$|$GM]}}
                \texttt{Answer for APPETIZER invalid after max attempts.} \\
            }
        }
    }
    & & \\ \\

    \theutterance \stepcounter{utterance}  
    & & & \multicolumn{2}{p{0.3\linewidth}}{
        \cellcolor[rgb]{0.9,0.9,0.9}{
            \makecell[{{p{\linewidth}}}]{
                \texttt{\tiny{[GM$|$GM]}}
                \texttt{Answer is incorrect.} \\
            }
        }
    }
    & & \\ \\

    \theutterance \stepcounter{utterance}  
    & & & \multicolumn{2}{p{0.3\linewidth}}{
        \cellcolor[rgb]{0.9,0.9,0.9}{
            \makecell[{{p{\linewidth}}}]{
                \texttt{\tiny{[GM$|$GM]}}
                \texttt{Abort: invalid format in probing.} \\
            }
        }
    }
    & & \\ \\

\end{supertabular}
}

\end{document}
